\documentclass[10pt,a4paper]{amsart}
\usepackage[margin=19mm]{geometry}
\usepackage{amsmath,amssymb,mathtools}
\usepackage[scr=rsfs,cal=euler]{mathalfa}
\usepackage{xfrac}
\usepackage[unicode,hidelinks,bookmarks=false]{hyperref}
\newcommand{\ie}{i.e.\ }
\newcommand{\cf}{cf.\ }
\newcommand{\resp}{respectively\ }
\newcommand{\Subsection}{\subsection}
\providecommand{\nobreakdash}{\nobreak\hbox{-}}
\setlength{\emergencystretch}{2em}
\multlinegap0pt
\usepackage{bm}
\usepackage{bbm}
\usepackage{dsfont}
\usepackage{xspace}
\usepackage{cancel}
\usepackage{mathtools}
\usepackage{amsthm}
\makeatletter
\@ifundefined{theorem}{\newtheorem{theorem}{Theorem}}{}
\@ifundefined{thm}{\newtheorem{thm}{Theorem}}{}
\makeatother
\newcommand{\C}{\mathbb{C}}
\title[Understanding Fermat's Last Theorem's Proofs]{Understanding Fermat's Last Theorem's Proofs}
\author{Alex Qiu, Tanish Sarathy, Spencer Nicklin, Michael Sun}
\date{Source: arXiv:2508.10362v1 (CC BY 4.0)}
\begin{document}
\maketitle

\noindent\emph{Source and licence.} Understanding Fermat's Last Theorem's Proofs. Version arXiv:2508.10362v1 (CC BY 4.0), distributed under the Creative Commons Attribution 4.0 International licence, as recorded in the arXiv licence metadata for this exact version and shown on the abstract page https://arxiv.org/abs/2508.10362v1. Full rights notice: licenses/arxiv-2508.10362-CC-BY-4.0.txt.\par\smallskip

\noindent\textit{Editorial scope.} Counted unit: the Thm whose source label is \texttt{None} in the article (source0.tex lines 2404--2415, source label \texttt{None}), delivered together with the article's own complete original proof of it (source0.tex lines 2417--2523, one \texttt{proof} environment of 107 lines). No mathematics is written, repaired or re-derived: statement and proof are the article's own text line for line. Required context retained uncounted: discretezero (lines 1856--1858); argprinciple (lines 1901--1904); extended-argument-principle (lines 4667--4671). Every definition, display and label of those lines is retained unchanged. Omitted from this package and not redistributed: 1--1855, 1859--1900, 1905--2403, 2416--2416, 2524--4666, 4672--5150. Nothing inside a retained excerpt is omitted or altered: every retained body is delivered exactly as the article's lines are written, so no omission receipt of any kind accompanies this package. Citations: the retained excerpts carry no citation command at all, so no bibliography entry of the article is transcribed and no reference is redistributed. Reference adaptation: none was needed and none was made; every reference inside the delivered unit resolves inside it, so no unresolved reference or citation is produced. Printed slips kept literal: no printed word, symbol, hypothesis or number of the counted statement or of its proof was altered. Layout and class: the article's own class is \texttt{article} with packages including graphicx, amsmath, amsfonts, amsthm, amssymb, xcolor, mathrsfs, breqn, comment, asymptote; the wrapper is the lane's pinned \texttt{amsart} editorial wrapper at 10pt on A4, which restates the theorem-like environments the retained text needs and adds the article's own macro definitions that the retained text needs (\texttt{\textbackslash C}), with extra wrapper packages (bm, bbm, dsfont, xspace, cancel, mathtools). The delivered numbering is the unit's own. Assets: no retained span contains a figure environment, a graphics-inclusion command, a PDF-page inclusion, a table or a TikZ picture; no image file is copied into this package, the package declares no asset dependency and no image file is redistributed with it. Licence: version arXiv:2508.10362v1 of this article is distributed under the Creative Commons Attribution 4.0 International licence, as recorded in the arXiv licence metadata for this exact version and shown on the abstract page https://arxiv.org/abs/2508.10362v1. A full rights notice is delivered at licenses/arxiv-2508.10362-CC-BY-4.0.txt. Characters: every retained excerpt is pure ASCII, so the delivered engine, XeLaTeX with its default Latin Modern text font, reports no missing character.
\begin{thm}\label{discretezero}
    The zeros of a holomorphic function are discrete. In other words, for any zero of a holomorphic function there exists a neighbourhood of that zero that contains no zero other than itself.
\end{thm}

\begin{thm}[Argument Principle]\label{argprinciple}For a meromorphic function $f$ and $C$ is a closed curve, then the zeros and poles enclosed by $C$ have difference given by
$$\frac{1}{2\pi i}\int_C\frac{f'(z)}{f(z)}dz= zeros-poles$$
counting multiplicity.
\end{thm}

\begin{theorem} \label{extended-argument-principle}
    If $U$ is a simply connected open subset of $\mathbb C$, $f$ is a holomorphic function on $U\setminus\{z_1,z_2,\dots,z_n\}$ where $z_1,z_2,\dots,z_n$ are its poles, and $\gamma$ is a simple rectifiable closed curve with winding number $1$ around every point of its interior, and there are a finite number of roots inside the curve, then the integral
    $$\int_\gamma \frac{zf'(z)}{f(z)}dz$$
    is equal to $2\pi i$ times the difference between the sum of the roots and the sum of the poles inside the curve, counting multiplicity.
\end{theorem}

\begin{thm} Let $\mathcal{L}$ be a lattice and $\C/\mathcal{L}$ the quotient as Abelian groups.
The equation

$$y^2=4x^3-60G_4(\mathcal{L})x-140G_6(\mathcal{L})$$
defines an elliptic curve over $\C$ (that is, it is non-singular) written $E_\C$.

Furthermore, there is an isomorphism of Abelian groups
$$\C/\mathcal{L}\to E_{\C}$$
given by
$$z+\mathcal{L}\mapsto (\wp(z),\wp'(z))$$
for $z\notin \mathcal{L}$ and $0\mapsto O$, the point at infinity.
\end{thm}

\begin{proof} Proof checklist: Bijective, Group Homomorphism.
\newline\newline
First we prove that the map creates a \textbf{bijection} between the Abelian groups. 
\newline\newline
Let $x\in\C$ and let $f(z)=\wp(z)-x$. We aim to show that $x$ is in the range of $\wp$ by showing $0$ is in the range of $f$ and with further examination of the zeros of $f$, conclude there is a bijection $z\mapsto (x,y)$.
\newline\newline
Due to our Laurent series expansion of $\wp$, we know that $f$ has a pole of multiplicity $2$ at $0$, so we can define $g(z)=z^2f(z)$ so then $g$ can be extended to be holomorphic at $0$. Now take a closed parallelogram of the lattice and shift it so that it is centered at $0$, that is,  let $\omega_1,\omega_2$ generate $\mathcal{L}$ and take the parallelogram with vertices $\pm\frac{\omega_1+\omega_2}{2}$ and $\pm\frac{\omega_1-\omega_2}{2}$ (including the boundary and interior points). So now the only lattice point on the parallelogram is $0$. Therefore $g$ is defined on the whole of the parallelogram. Call this parallelogram $\mathcal{P}$.
\newline\newline
We show that $f$ has only finitely many zeros in $\mathcal{P}$.
Assume for the sake of contradiction that $f$ had infinitely many zeroes in $\mathcal{P}$. Then $g$ has infinitely many zeroes in $\mathcal{P}$, and we can form an infinite sequence of distinct zeroes in $\mathcal{P}$, so by the Bolzano-Weierstrass theorem, there is a convergent subsequence, so this sequence has a limit point inside the closed parallelogram. Then $g=0$ so $f=0$ in the parallelogram (Theorem \ref{discretezero}), so due to periodicity then $f=0$ everywhere, which is clearly not true, looking at the Laurent series of $f$ and we have reached a contradiction.
\newline\newline
Now since there are a finite number of zeroes on any period of the lattice, shift the parallelogram so that it has no zeroes nor lattice points on its border. Call this shifted parallelogram $\mathcal{P}_1$ and let this parallelogram's counterclockwise boundary be $\partial \mathcal{P}_1$. We show that
$$\int_{\partial \mathcal{P}_1} \frac{f'(z)}{f(z)} dz=0$$
and use the argument principle (Theorem \ref{argprinciple}), to count the number of zeros of $f$.
\newline\newline
Due to $\mathcal L$-periodicity of $f'/f$, the opposite sides of the parallelogram being in opposite directions cancel out in the integral, so this evaluates to $0$. Thus using the argument principle, $f$ has the same number of zeroes on the parallelogram as poles counting multiplicity, and there are $2$ poles (the double pole at the lattice point), so it has $2$ zeroes counting multiplicity. Therefore, $\wp(z)=x$ for either two distinct values of $z$ or a single value with multiplicity two, up to translations of $\mathcal{L}$ (we will say modulo $\mathcal{L}$).
\newline\newline
We now show that for all but three $z$ values modulo $\mathcal{L}$ (keep in mind that a cubic has three roots), there are exactly two distinct values of $z$ that produce distinct $y$ values (ie non-zero $y^2$) up to translation by $\mathcal{L}$, otherwise, $y^2=0$.
\newline\newline
For any $z\in\mathbb C\setminus\mathcal L$ not equal to $\frac{\omega_1}2$, $\frac{\omega_2}2$ or $\frac{\omega_1+\omega_2}2$ $\mod\mathcal{L}$, $z\not\equiv-z\mod\mathcal L$, so since $\wp$ is even, $z,-z$ are those values mod $\mathcal{L}$ on which $\wp(-z)$ equals $\wp(z)$. Hence in this case, there are exactly two distinct $z$ that output the same $x$. We show these output non-zero $y$-values which are negatives of each other. To see they are negatives of each other, we recall that $\wp'(-z)=-\wp'(z)$ is odd. To show these $\pm y$ are non-zero, it suffices to show $\wp(\frac{\omega_1}2)$, $\wp(\frac{\omega_2}2)$ and $\wp(\frac{\omega_1+\omega_2}2)$ are distinct zeros different from $\wp(z)$ since a cubic equation in $x$ can have at most 3 zeros.
\newline\newline
Since $\wp'$ is odd, that means that $$\wp'(\frac{\omega_1}2)=-\wp'(-\frac{\omega_1}2)=-\wp'(-\frac{\omega_1}2+\omega_1)=-\wp'(\frac{\omega_1}2),$$
so 
$$\wp'(\frac{\omega_1}2)=0.$$ Similarly, $$\wp'(\frac{\omega_2}2)=0\qquad \text{and}\qquad \wp'(\frac{\omega_1+\omega_2}2)=0.$$ 
We now explain why these are different to the other $\wp(z)$ using multiplicity of zeros of $f$ established above, namely, we show these each have multiplicity 2. If $\alpha\in\{\frac{\omega_1}2,\frac{\omega_2}2,\frac{\omega_1+\omega_2}2\}$, then $f(\alpha)=0$ for $f$ with $x=\wp(\alpha)$ and $\alpha$ has multiplicity at least two since $f'(\alpha)=\wp'(\alpha)=0$. This means that $\alpha$ has multiplicity exactly 2 using the multiplicity of zeros of $f$ already established.
\newline\newline
% then these three must each have multiplicity at least two and hence must be equal to two. 

%This show

%If everything in $\{\frac{\omega_1}2,\frac{\omega_2}2,\frac{\omega_1+\omega_2}2\}$ has multiplicity 2, then we are done. Otherwise, three different $\alpha,\beta,\gamma\in\{\frac{\omega_1}2,\frac{\omega_2}2,\frac{\omega_1+\omega_2}2\}$ satisfy $\wp(\alpha)=\wp(\beta)$ and $\gamma$ has multiplicity 2. This is because using other values $z$ will lead to three zeros of $f$.

%Therefore, $\wp$ takes those values with multiplicity at least $2$, and hence those points are the unique points mod the lattice with their particular $\wp$ values. 
This also shows that $\wp(\frac{\omega_1}2),\wp(\frac{\omega_2}2),\wp(\frac{\omega_1+\omega_2}2)$ are distinct roots of the polynomial $4x^3-60G_4(\mathcal L)x-140G_6(\mathcal L)$. Hence the curve $y^2=4x^3-60G_4(\mathcal L)x-140G_6(\mathcal L)$ is non-singular. 
\newline\newline
%, with the roots of the cubic being precisely those values, and no other points mod the lattice can be roots of $\wp'$.
In summary, $\frac{\omega_1}2,\frac{\omega_2}2,\frac{\omega_1+\omega_2}2$ are taken to the three distinct $x$-values which are roots of the cubic, at which the elliptic curve only assumes one $y$-value of zero, and all other points modulo $\mathcal{L}$ are taken to $x$-values at which the elliptic curve assumes two $y$-values which are negatives of each other, and each $y$-value is taken once due to $\wp'$ being odd. Finally if we extend the map to take the lattice points to the point at infinity on the elliptic curve, this map bijects $\mathbb C/\mathcal L$ with the elliptic curve.
\newline\newline
Now we need to show that this map is a \textbf{group homomorphism}, in other words, it is additive. Say that we have two points, $z_1$ and $z_2$, and we want to show that $z_1+z_2$ gets taken to the sum of the corresponding points on the elliptic curve. First we deal with the case where one of them is a lattice point. If $z_1\in\mathcal L$ then $z_1+z_2\equiv z_2\mod\mathcal L$, so $z_1+z_2$ gets taken to the same point as $z_2$ as expected. Otherwise, $z_1,z_2$ get taken to actual points in $\mathbb C^2$ by $\wp$ and $\wp'$. If $(\wp(z_1),\wp'(z_1))$ and $(\wp(z_2),\wp'(z_2))$ are inverses, then $\wp(z_1)=\wp(z_2)$ and $\wp'(z_1)=-\wp'(z_2)\neq\infty$, it can be checked from the bijection above that $z_1+z_2\equiv0\mod\mathcal L$, and hence $z_1+z_2$ gets sent to the point at infinity as expected. Otherwise, the line joining the points $(\wp(z_1),\wp'(z_1))$ and $(\wp(z_2),\wp'(z_2))$ isn't "vertical" since otherwise it would pass through the point at infinity (Note: this includes the tangent case as well), so it can be expressed as $y=\ell(x)$ where $\ell$ is a linear polynomial. Define
$$h(z)=\ell(\wp(z))-\wp'(z)$$
which then has roots corresponding exactly to the intersections of the line and the elliptic curve. Define 
$$p(x)=4x^3+60G_4(\mathcal L)x-140G_6(\mathcal L).$$ Given such a root $w$ of $h$, then we know that $\wp'(w)=\ell(\wp(w))$ since the point is on the line. The differential equation for $\wp$ can be expressed as $$\wp'(z)^2=p(\wp(z)),$$
so then differentiating this we get $2\wp'(z)\wp''(z)=\wp'(z)p'(\wp(z))$. Therefore for all $z$ where $\wp'(z)\neq0$, 
$$2\wp''(z)=p'(\wp(z)),$$
and since the zeroes of a non-zero holomorphic function ($\wp'$) are isolated, by continuity we get this identity for all $z\in\mathbb C\setminus\mathcal L$. Therefore, given that $w$ is a root of $h$, then $\wp'(w)^2=\ell(\wp(w))^2=p(\wp(w))$, so $\wp(w)$ is a root of $\ell^2-p$, then
\begin{align*}
&\text{$w$ is at least a double root of $h$}\\
   \iff&h'(w)&=0\\ 
   \iff& \wp'(w)\ell'(\wp(w))-\wp''(w)&=0\\ 
   \iff&2\ell(\wp(w))\ell'(\wp(w))-p'(\wp(w))&=0\\
   \iff&(\ell^2-p)'(\wp(w))&=0\\
   \iff&\text{$\wp(w)$ is at least a double root of $\ell^2-p$}
\end{align*}
Now assume further that the above equivalence is true for $w$. Then we can show that $\ell(\wp(w))\neq0$: assume for the sake of contradiction that $\ell(\wp(w))=0$. then $\wp'(w)=\ell(\wp(w))=0$ and $\ell(\wp(w))^2=p(\wp(w))=0$, and also from the above equivalence (second $\iff$), we have $\wp'(w)\ell'(\wp(w))-\wp''(w)=0$, then $\wp''(w)=\wp'(w)\ell'(\wp(w))=0$. Finally, $p'(\wp(w))=2\wp''(w)=0$. Now we have $p(\wp(w))=0$ and $p'(\wp(w))=0$, so this is a contradiction since that means $p$ has a double root at $\wp(w)$, but we already established above that $p$ has $3$ distinct roots. 

Therefore since $\ell(\wp(w))\neq0$, from $2\ell(\wp(w))\ell'(\wp(w))-p'(\wp(w))=0$, we can conclude that 
$$\ell'(\wp(w))=0\iff p'(\wp(w))=0.$$
Furthermore, differentiating the identity $2\wp''(z)=p'(\wp(z))$ we get $$2\wp'''(z)=\wp'(z)p''(\wp(z)).$$
Also since $\ell$ is linear, $\ell''=0$. Hence
\begin{align*}
    &\text{$w$ is at least a triple root of $h$}\\
    \iff &h''(w)&=0\\
    \iff &\wp''(w)\ell'(\wp(w))+\wp'(w)^2\ell''(\wp(w))-\wp'''(w)&=0\\
    \iff &p'(\wp(w))\ell'(\wp(w))-\wp'(w)p''(\wp(w))&=0\\
    \iff &p'(\wp(w))\ell'(\wp(w))-\ell(\wp(w))p''(\wp(w))&=0
\end{align*}

We also have that
\begin{align*}
    &\text{$\wp(w)$ is at least a triple root of $\ell^2-p$}\\
    \iff &(\ell^2-p)''(\wp(w))&=0\\
    \iff &(2{\ell'}^2+2\ell\ell''-p'')(\wp(w))&=0\\
    \iff &(2{\ell'}^2-p'')(\wp(w))&=0\\
    \iff&2\ell'(\wp(w))^2-p''(\wp(w))&=0
\end{align*}
In the case that $\ell'(\wp(w))=p'(\wp(w))=0$, then 
\begin{align*}
&2\ell'(\wp(w))^2-p''(\wp(w))&=0\\
\iff&p''(\wp(w))&=0\\
\iff&\ell(\wp(w))p''(\wp(w))&=0\\
\iff&p'(\wp(w))\ell'(\wp(w))-\ell(\wp(w))p''(\wp(w))&=0
\end{align*}
In the case that $\ell'(\wp(w))\neq0$ and $p'(\wp(w))\neq0$, then
\begin{align*}
&2\ell'(\wp(w))^2-p''(\wp(w))&=0\\
\iff&2\ell'(\wp(w))^2p'(\wp(w))-p''(\wp(w))p'(\wp(w))&=0\\
\iff&\ell'(\wp(w))^2p'(\wp(w))-p''(\wp(w))\ell(\wp(w))\ell'(\wp(w))&=0\\
\iff&\ell'(\wp(w))p'(\wp(w))-p''(\wp(w))\ell(\wp(w))&=0\\
\end{align*}
so either way, 
$$w\text{ is at least a triple root of }h \iff \wp(w)\text{ is at least a triple root of }\ell^2-p$$
$\ell^2-p$ has degree $3$ so roots can be at most triple roots. As for $h$, looking at the definition, $h$ has poles of multiplicity $3$ at each lattice point (that is from the $-\wp'$ term), so applying a similar argument that we did to $\wp$ above, except $g(z)=z^3h(z)$ this time, we get that $h$ has $3$ zeroes counting multiplicity on a period of the lattice. Hence roots of $h$ also can have multiplicity at most $3$. Therefore,  the multiplicity of the root $w$ of $h$ is equivalent to the multiplicity of the root $\wp(w)$ of $\ell^2-p$ where $\wp(w)$ is the $x$-coordinate of the intersection, i.e. the multiplicity of the intersection of the line with the elliptic curve.
\newline\newline
Now this means that mod $\mathcal L$, $z_1$ and $z_2$ are two of the roots of $h$ counting multiplicity, and the third root $z_3$ counting multiplicity corresponds to the third intersection of the line with the elliptic curve counting multiplicity. Now consider once again the anticlockwise border of a parallelogram of the lattice $\gamma=\gamma_1+\gamma_2+\gamma_3+\gamma_4$ where $\gamma_1,\gamma_2,\gamma_3,\gamma_4$ are its sides, $\gamma_3$ is $\gamma_1$ translated by $\omega_1$ and reversed, $\gamma_4$ is $\gamma_2$ translated by $\omega_2$ and reversed, and $\gamma$ doesn't pass through any poles or zeroes of $h$. Then consider
$$\int_\gamma \frac{zh'(z)}{h(z)}dz$$
which turns out from Theorem \ref{extended-argument-principle} to equal $2\pi i$ times the difference between the sum of the roots $z_1+z_2+z_3$ and the sum of poles of $h$ inside the parallelogram, counting multiplicity. We now show this integral evaluates to $2\pi i$ times an element of $\mathcal{L}$:
{\tiny
\begin{align*}
\int_\gamma \frac{zh'(z)}{h(z)}dz&=\int_{\gamma_1} \frac{zh'(z)}{h(z)}dz+\int_{\gamma_2} \frac{zh'(z)}{h(z)}dz+\int_{\gamma_3} \frac{zh'(z)}{h(z)}dz+\int_{\gamma_4} \frac{zh'(z)}{h(z)}dz\\
&=\int_{\gamma_1} \frac{zh'(z)}{h(z)}dz+\int_{\gamma_2} \frac{zh'(z)}{h(z)}dz-\int_{\gamma_1} \frac{(z+\omega_1)h'(z+\omega_1)}{h(z+\omega_1)}dz-\int_{\gamma_2} \frac{(z+\omega_2)h'(z+\omega_2)}{h(z+\omega_2)}dz\\
&=\int_{\gamma_1} \frac{zh'(z)}{h(z)}dz+\int_{\gamma_2} \frac{zh'(z)}{h(z)}dz-\int_{\gamma_1} \frac{(z+\omega_1)h'(z)}{h(z)}dz-\int_{\gamma_2} \frac{(z+\omega_2)h'(z)}{h(z)}dz\\
&=-\int_{\gamma_1} \frac{\omega_1h'(z)}{h(z)}dz-\int_{\gamma_2} \frac{\omega_2h'(z)}{h(z)}dz\\
&=-\omega_1\int_{\gamma_1} \frac{h'(z)}{h(z)}dz-\omega_2\int_{\gamma_2} \frac{h'(z)}{h(z)}dz.
\end{align*}
}
From here it suffices to show the remaining integrals are integers. These are two integrals of the logarithmic derivative of $h$ starting and ending at points which, due to the $\mathcal L$-periodicity of $h$, have equal values of $h$. Therefore each integral evaluates to an integer multiple of $2\pi i$, making the above evaluate to $2\pi i$ times a lattice point. Also, the only pole inside the parallelogram is a triple pole at a lattice point. Therefore the sum of the roots of $h$ in the parallelogram is itself actually part of the lattice, i.e. $z_1+z_2+z_3\in\mathcal L$ and hence $z_1+z_2+z_3\equiv0\mod\mathcal L$. Then since $z_3\equiv-z_1-z_2\mod\mathcal L$ and maps to the third intersection of the line with the elliptic curve, i.e. the negative of the sum of the other two points on the elliptic curve, then $-z_3\equiv z_1+z_2\mod\mathcal L$ maps to the actual sum of the other two points on the elliptic curve. Hence the map is a group isomorphism.
\end{proof}

\end{document}
